% Part: proof-theory
% Chapter: normalization
% Section: introduction

\documentclass[../../../include/open-logic-section]{subfiles}

\begin{document}

\olfileid{pt}{nor}{int}

\olsection{Pambuka}

Yèn sampeyan wis !!{prove} implikasi~$!A \lif !B$, bisa waé
sampeyan migunakaké ing !!a{proof} kanggo~$!B$ kanthi
\Elim{\lif}. Mesthiné iki uga mbutuhaké !!a{proof}~$\delta_1$
kanggo~$!A$. !!{proof} sampeyan kanggo~$!A \lif !B$
lumrahé pungkasané \Intro{\lif}, kang ngowahi
!!a{proof}~$\delta_2$ kanggo~$!B$ saka asumsi kang isih
mbukak~$!A$ dadi bukti kanggo~$!A \lif !B$. Yèn mangkono,
!!{proof} pungkasané wujudé
\[
\AxiomC{$\Discharge{!A}{x}$}
\RightLabel{$\delta_2$}
\DeduceC{$!B$}
\DischargeRule{\Intro{\lif}}{x}
\UnaryInfC{$!A \lif !B$}
\AxiomC{}
\RightLabel{$\delta_1$}
\DeduceC{$!A$}
\RightLabel{\Elim{\lif}}
\BinaryInfC{$!B$}
\DisplayProof
\]
Senajan siasat sampeyan irit amarga migunakaké manèh apa kang wis
ana (!!a{proof} kanggo~$!A \lif !B$), !!{proof}é dhéwé ora irit.
Kaya ana liku-liku nalika $!A \lif !B$ dilebokaké mung kanggo
langsung diilangaké. Kuduné ana cara kang luwih langsung kanggo
!!{prove}~$!B$.

Pancen, cara kuwi mesthi ana. ``Liku-liku'' kaya ing ndhuwur---aturan
!!a{introduction} kang langsung disusul aturan
!!a{elimination}---mesthi bisa diilangi saka !!{proof}s deduksi alami.
Iki diarani \emph{normalisasi}, lan asilé yaiku !!{proof}
\emph{normal} (utawa !!a{proof} \emph{ing wujud normal}).

Mbuktèkaké asil iki, kaya teorema penghilangan cut ing kalkulus
sekuen, rada ruwet lan mbutuhaké pamriksan akèh kasus. Kanggo
\Log{N1c} klasik, buktiné luwih angel tinimbang kanggo
\Log{N1i} intuisionistik. Kaya déné penghilangan cut nuduhaké yèn
aturan \Cut{} ing kalkulus sekuen ora ndadèkaké luwih akèh perkara
bisa !!{prove}, normalisasi nuduhaké yèn liku-liku ora ndadèkaké
luwih akèh perkara bisa dibuktèkaké tinimbang nalika liku-liku mau
diendhani. Ana pepadhan liyané: !!{proof} normal bisa luwih dawa
tinimbang !!{proof} paling cekak kang mawa liku-liku kanggo asil
kang padha, kaya déné !!{proof}s ing kalkulus sekuen bisa luwih
cekak kanthi cut tinimbang !!{proof} tanpa cut kang paling cekak.
Ing deduksi alami intuisionistik, sipat sub-!!{formula}, yaiku yèn
sampeyan tansah bisa !!{prove} sawijining asil nganggo mung
sub-!!{formula}s saka asumsi kang isih mbukak utawa saka asilé ing
!!{proof} sampeyan, ndhèrèk saka normalisasi, kaya déné sipat iku
ndhèrèk saka penghilangan cut kanggo kalkulus sekuen.

\end{document}
